\documentclass{article}
\usepackage[utf8]{inputenc}
\usepackage[T1]{fontenc}
\usepackage[french]{babel}
\usepackage{amsmath}
\usepackage{amsfonts}
\usepackage{amssymb}
\usepackage{xcolor}
\usepackage{tikz}
\usepackage{titlesec}

% --- Configuration automatique des couleurs pour les titres ---
% Change la couleur des sections (\section) en rouge
\titleformat{\section}
{\color{red}\normalfont\Large\bfseries}{\thesection}{1em}{\underline}

% Change la couleur des sous-sections (\subsection) en rouge
\titleformat{\subsection}
{\color{red}\normalfont\large\bfseries}{\thesubsection}{1em}{\underline}

% Change la couleur des sous-sous-sections (\subsubsection) en rouge
\titleformat{\subsubsection}
{\color{red}\normalfont\normalsize\bfseries}{\thesubsubsection}{1em}{\underline}

\begin{document}

\section*{Chap : Barycentre}

\subsection*{Introduction}

\noindent Sachant que la balance ci-dessous est en équilibre au point $G$.

\vspace{1em}

\begin{center}
    \begin{tikzpicture}[thick, scale=0.8]
        % La barre de la balance
        \draw (0,0) -- (6,0);
        
        % Les masses (boîtes)
        \draw (0,0) rectangle (1,1) node[midway, above=0.5cm] {$m_A$};
        \draw (5,0) rectangle (6,1) node[midway, above=0.5cm] {$m_B$};
        
        % Les points
        \node at (0.5,-0.3) {$A$};
        \node at (5.5,-0.3) {$B$};
        \node at (3,0) {$\times$};
        \node at (3,-0.3) {$G$};
    \end{tikzpicture}
\end{center}

\vspace{1em}

\noindent Où $m_A$ et $m_B$ sont les masses respectives des objets posés au point $A$ et $B$. \\
Alors, d'après les lois de la physique, on a : \\
$m_A \overrightarrow{AG} = m_B \overrightarrow{GB}$ \\
$-m_A \overrightarrow{GA} = m_B \overrightarrow{GB}$ \\
$-m_A \overrightarrow{GA} - m_B \overrightarrow{GB} = \vec{0}$ \\
$m_A \overrightarrow{GA} + m_B \overrightarrow{GB} = \vec{0}$

\vspace{1em}

\noindent En physique, on dit que $G$ est le \textbf{centre d'inertie}. \\
En mathématiques, on dit que $G$ est le \textbf{barycentre}.
\section*{I- Barycentre de deux points pondérés}

\subsection*{1- Notions de base}

\subsubsection*{Définition}

\noindent On appelle \textbf{point pondéré} tout couple $(P, m)$ où $P$ est un point du plan et $m$ un réel appelé \textbf{coefficient} de $P$ ou \textbf{masse} de $P$.

\vspace{1em}

\noindent \underline{\textbf{Ex :}} \\
$(A, 2)$ est un point pondéré. \\
$(B, -3)$ est un point pondéré.

\subsection*{Théorème - Définition}

\noindent Soient $(A, a)$ et $(B, b)$ deux points pondérés. \\
Si $a + b \neq 0$, alors il existe un unique point $G$ du plan tel que :
\[ a \overrightarrow{GA} + b \overrightarrow{GB} = \vec{0} \]
\noindent Le point $G$ est appelé \textbf{barycentre} des points pondérés $(A, a)$ et $(B, b)$, on note :

\vspace{1em}

\noindent $G = \text{bar} \{(A, a), (B, b)\}$

\vspace{1em}

\noindent $G = \text{bar} \begin{tabular}{|c|c|}
\hline
$A$ & $B$ \\
\hline
$a$ & $b$ \\
\hline
\end{tabular}$

\subsection*{Preuve :}

\noindent On a : $a \overrightarrow{GA} + b \overrightarrow{GB} = \vec{0}$ \\
$a \overrightarrow{GA} + b (\overrightarrow{GA} + \overrightarrow{AB}) = \vec{0}$ \\
$a \overrightarrow{GA} + b \overrightarrow{GA} + b \overrightarrow{AB} = \vec{0}$ \\
$(a + b) \overrightarrow{GA} + b \overrightarrow{AB} = \vec{0}$ \\
$-(a + b) \overrightarrow{AG} + b \overrightarrow{AB} = \vec{0}$ \\
$-(a + b) \overrightarrow{AG} = -b \overrightarrow{AB}$ \\
$(a + b) \overrightarrow{AG} = b \overrightarrow{AB}$ \\
$\overrightarrow{AG} = \dfrac{b}{(a + b)} \overrightarrow{AB}$

\vspace{1.5em}

\noindent En posant $\vec{u} = \dfrac{b}{a + b} \overrightarrow{AB}$ \\
On obtient $\overrightarrow{AG} = \vec{u}$

\vspace{1em}

\noindent Donc d'après la propriété fondamentale, $G$ est l'unique point tel que

 $a \overrightarrow{GA} + b \overrightarrow{GB} = \vec{0}$ avec $a + b \neq 0$.
 
 \subsection*{Remarque :}

\noindent \textbf{1) Compte tenu de la preuve du théorème - définition, on a :}

\vspace{0.5em}

\begin{itemize}
    \item[i)] $G = \text{bar} \{(A, a), (B, b)\} \iff \overrightarrow{AG} = \dfrac{b}{a+b} \overrightarrow{AB}$
    \item[ii)] $G = \text{bar} \{(A, a), (B, b)\} \iff \overrightarrow{BG} = \dfrac{a}{a+b} \overrightarrow{BA}$
\end{itemize}

\vspace{0.5em}

\noindent \textbf{Donc pour construire le point $G$, on applique i) ou ii).}

\vspace{1em}

\noindent \textbf{\color{red}2) Si $a = b$ alors $A = B = G$} (Note : Cas de confusion de points dans le plan)

\vspace{0.5em}

\noindent \textbf{\color{red}3) Si $G = \text{bar} \{(A, a), (B, b)\}$ alors $G \in (AB)$}

\subsection*{Application}

\noindent Soit $G = \text{bar} \{(A, 2), (B, 1)\}$ \\
Construire le point $G$.

\vspace{1.5em}

\begin{center}
    \underline{\color{red}Correction}
\end{center}

\noindent $\overrightarrow{AG} = \dfrac{b}{a + b} \overrightarrow{AB} \implies \overrightarrow{AG} = \dfrac{1}{2 + 1} \overrightarrow{AB} \implies \overrightarrow{AG} = \dfrac{1}{3} \overrightarrow{AB}$

\vspace{2em}

\begin{center}
    \begin{tikzpicture}[thick, scale=1.5]
        % Tracé du segment principal
        \draw (0,0) -- (4,0);
        
        % Graduations (tiers du segment)
        \foreach \x in {0, 1.33, 2.66, 4}
            \draw (\x, 0.1) -- (\x, -0.1);
            
        % Marquage des points
        \node[below] at (0,-0.1) {$A$};
        \node[below] at (4,-0.1) {$B$};
        \node at (1.33, 0) {$\times$};
        \node[below] at (1.33,-0.1) {$G$};
        
        % Triangle de construction manuscrit
        \draw (0,0) -- (2,-1) -- (4,0);
    \end{tikzpicture}
\end{center}

\subsubsection*{Définition :}

Le barycentre des points pondérés affecté des coefficients égaux est appelé \textbf{isobarycentre}.

\vspace{1em}

\subsubsection*{Exemple :}

\noindent \textbf{i)} $G = \text{bar} \{(A, 4), (B, 4)\}$ \\
\textbf{ii)} $G = \text{bar} \{(A, -3), (B, -3)\}$

\vspace{1em}

\subsubsection*{Remarque :}

L'isobarycentre de deux points $A$ et $B$ distincts est le \textbf{milieu} du segment $[AB]$.

\vspace{1em}

\noindent \textit{En effet,} \\
Soit $G = \text{bar} \{(A, b), (B, a)\}$ \\
$G$ est l'isobarycentre des points $A$ et $B$ lorsque $a = b$. \\
$G = \text{bar} \{(A, b), (B, b)\} \implies b \overrightarrow{GA} + b \overrightarrow{GB} = \vec{0}$ \\
\phantom{$G = \text{bar} \{(A, b), (B, b)\} \implies$} or $a = b$ donc $b \overrightarrow{GA} + b \overrightarrow{GB} = \vec{0}$ \\
\phantom{$G = \text{bar} \{(A, b), (B, b)\} \implies$} $b (\overrightarrow{GA} + \overrightarrow{GB}) = \vec{0}$ \\
\phantom{$G = \text{bar} \{(A, b), (B, b)\} \implies$} Donc $\overrightarrow{GA} + \overrightarrow{GB} = \vec{0}$ \\
\phantom{$G = \text{bar} \{(A, b), (B, b)\} \implies$} $\overrightarrow{GA} = -\overrightarrow{GB}$ \\
\textbf{D'où $G$ est le milieu de $[AB]$.}

\subsubsection*{Exemple :}

$G = \text{bar} \{(A, 1), (B, 1)\}$ \\
donc $G$ est le milieu de $[AB]$.

\vspace{2em}

\subsection*{2- Propriétés}

\subsubsection*{Propriété 1 : Homogénéité du Barycentre}

Le barycentre de deux points pondérés est inchangé lorsqu'on multiplie tous les points par un même nombre réel non nul.

\vspace{1em}

$G = \text{bar} \{(A, a), (B, b)\} \iff G = \text{bar} \{(A, ka), (B, kb)\}$ avec $k \neq 0$

\subsubsection*{Exemple :}

\noindent 1) $G = \text{bar } \{(A, 1), (B, 2)\}$ \\
Si $K = 5$ \\
$G = \text{bar } \{(A, 5), (B, 10)\}$ \\
Si $K = -2$ \\
$G = \text{bar } \{(A, -2), (B, -4)\}$

\vspace{1em}

\noindent 2) $G = \text{bar } \{(A, 4), (B, 6)\}$ \\
Si $K = \dfrac{1}{2}$ \\
$G = \text{bar } \{(A, 2), (B, 3)\}$

\vspace{2em}

\subsection*{Propriété 2 : Réduction de la somme $a\overrightarrow{MA} + b\overrightarrow{MB}$}

\noindent Soient $(A, a)$ et $(B, b)$ deux points pondérés tel que $a + b \neq 0$. \\
Pour tout point $M$ du plan, on a : \\
$G = \text{bar } \{(A, a), (B, b)\} \iff a\overrightarrow{MA} + b\overrightarrow{MB} = (a + b)\overrightarrow{MG}$

\subsection*{Preuve :}

\noindent $G = \text{bar } \{(A, a), (B, b)\} \iff a \overrightarrow{GA} + b \overrightarrow{GB} = \vec{0}$ \\
\phantom{$G = \text{bar } \{(A, a), (B, b)\} \iff$} $a (\overrightarrow{GM} + \overrightarrow{MA}) + b (\overrightarrow{GM} + \overrightarrow{MB}) = \vec{0}$ \\
\phantom{$G = \text{bar } \{(A, a), (B, b)\} \iff$} $a \overrightarrow{GM} + a \overrightarrow{MA} + b \overrightarrow{GM} + b \overrightarrow{MB} = \vec{0}$ \\
\phantom{$G = \text{bar } \{(A, a), (B, b)\} \iff$} $(a + b) \overrightarrow{GM} + a \overrightarrow{MA} + b \overrightarrow{MB} = \vec{0}$ \\
\phantom{$G = \text{bar } \{(A, a), (B, b)\} \iff$} $a \overrightarrow{MA} + b \overrightarrow{MB} = -(a + b) \overrightarrow{GM}$ \\
\phantom{$G = \text{bar } \{(A, a), (B, b)\} \iff$} $a \overrightarrow{MA} + b \overrightarrow{MB} = (a + b) \overrightarrow{MG}$

\vspace{2em}

\subsection*{Remarque 1 :}

\noindent En vertu de la preuve de la propriété 2), pour tout point $M$ du plan, on a : \\
$G = \text{bar } \{(A, a), (B, b)\} \iff \overrightarrow{MG} = \dfrac{1}{(a + b)} \left( a \overrightarrow{MA} + b \overrightarrow{MB} \right)$

\subsection*{Remarque 2 :}

Si $a + b = 0$ alors partout pour tout point $M$, le vecteur $a\overrightarrow{MA} + b\overrightarrow{MB}$ est indépendant du point $M$.

\vspace{1em}

\noindent \textit{En effet,}
\begin{align*}
    a\overrightarrow{MA} + b\overrightarrow{MB} &= a\overrightarrow{MA} + b(\overrightarrow{MA} + \overrightarrow{AB}) \\
    &= a\overrightarrow{MA} + b\overrightarrow{MA} + b\overrightarrow{AB} \\
    &= (a + b)\overrightarrow{MA} + b\overrightarrow{AB}
\end{align*}

\noindent Donc si $a + b = 0$ alors 
\begin{center}
    \color{red}
    \framebox{\color{black} $a\overrightarrow{MA} + b\overrightarrow{MB} = b\overrightarrow{AB}$}
\end{center}

\noindent D'où $a\overrightarrow{MA} + b\overrightarrow{MB}$ est indépendant de $M$.

\section*{NB :}

\noindent Soient $A$ et $B$ deux points distincts.

\vspace{1em}

\noindent \textbf{*1) Si $MA = MB$ alors l'ensemble des points $M$ est la médiatrice du segment $[AB]$.}

\vspace{1.5em}

\noindent Soient $R > 0$ et $G$ un point.

\vspace{1em}

\noindent \textbf{*2) Si $MG = R$ alors l'ensemble des points $M$ est un cercle de centre $G$ et de rayon $R$.}

\section*{Application}

\noindent $AB = 5\text{ cm}$

\vspace{1em}

\noindent \textbf{1) Soit $G$ barycentre des points pondérés $\{(A, 2), (B, 3)\}$} \\
Montrer que pour tout point $M$ : \\
$2 \overrightarrow{MA} + 3 \overrightarrow{MB} = 5 \overrightarrow{MG}$

\vspace{1.5em}

\noindent \textbf{2) Montrer que le vecteur} \\
$3 \overrightarrow{MA} - 3 \overrightarrow{MB}$ est indépendant du point $M$.

\vspace{1.5em}

\noindent \textbf{3) Déterminer l'ensemble des points $M$ du plan dans les cas suivants :}

\vspace{1em}

\begin{itemize}
    \item[a)] $\| 2 \overrightarrow{MA} + 3 \overrightarrow{MB} \| = \| 3 \overrightarrow{MA} - 3 \overrightarrow{MB} \|$
    \item[b)] $\| 2 \overrightarrow{MA} + 3 \overrightarrow{MB} \| = \| 3 \overrightarrow{MA} + 2 \overrightarrow{MB} \|$
\end{itemize}

\section*{Solution}

\noindent \textbf{1) Montrons que pour tout point $M$ :} \\
$2 \overrightarrow{MA} + 3 \overrightarrow{MB} = 5 \overrightarrow{MG}$

\vspace{1em}

\begin{align*}
    2 \overrightarrow{MA} + 3 \overrightarrow{MB} &= 2 (\overrightarrow{MG} + \overrightarrow{GA}) + 3 (\overrightarrow{MG} + \overrightarrow{GB}) \\
    &= 2 \overrightarrow{MG} + 2 \overrightarrow{GA} + 3 \overrightarrow{MG} + 3 \overrightarrow{GB} \\
    &= 5 \overrightarrow{MG} + 2 \overrightarrow{GA} + 3 \overrightarrow{GB}
\end{align*}

\vspace{1.5em}

\noindent or $G = \text{bar } \{(A, 2), (B, 3)\} \implies 2 \overrightarrow{GA} + 3 \overrightarrow{GB} = \vec{0}$ \\
D'où 
\begin{center}
    \framebox{$2 \overrightarrow{MA} + 3 \overrightarrow{MB} = 5 \overrightarrow{MG}$}
\end{center}

\noindent \underline{Méthode 2 :} $a = 2$ et $b = 3$ \\
$a \overrightarrow{MA} + b \overrightarrow{MB} = (a + b) \overrightarrow{MG}$ \\
$2 \overrightarrow{MA} + 3 \overrightarrow{MB} = (2 + 3) \overrightarrow{MG}$ \\
$2 \overrightarrow{MA} + 3 \overrightarrow{MB} = 5 \overrightarrow{MG}$

\vspace{2em}

\noindent \textbf{2) Montrons que le vecteur $3 \overrightarrow{MA} - 3 \overrightarrow{MB}$ est indépendant du point $M$.}

\vspace{1em}

\begin{align*}
    3 \overrightarrow{MA} - 3 \overrightarrow{MB} &= 3 \overrightarrow{MA} - 3 (\overrightarrow{MA} + \overrightarrow{AB}) \\
    &= 3 \overrightarrow{MA} - 3 \overrightarrow{MA} - 3 \overrightarrow{AB} \\
    &= -3 \overrightarrow{AB}
\end{align*}

\vspace{1em}

\noindent $3 \overrightarrow{MA} - 3 \overrightarrow{MB} = -3 \overrightarrow{AB}$ \\
\textbf{Donc le vecteur $3 \overrightarrow{MA} - 3 \overrightarrow{MB}$ est indépendant du point $M$.}

\noindent \underline{Méthode 2 :} \\
On a : $3 + (-3) = 0$ donc le vecteur $3\overrightarrow{MA} - 3\overrightarrow{MB}$ est indépendant du point $M$.

\vspace{1.5em}

\noindent \textbf{3) Déterminons l'ensemble des points $M$ dans les cas suivants :}

\vspace{1em}

\noindent \textbf{a)} $\| 2\overrightarrow{MA} + 3\overrightarrow{MB} \| = \| 3\overrightarrow{MA} - 3\overrightarrow{MB} \|$

\vspace{1em}

\noindent D'après ce qui précède on a : \\
$2\overrightarrow{MA} + 3\overrightarrow{MB} = 5\overrightarrow{MG}$ avec $G = \text{bar } \{(A, 2), (B, 3)\}$ \\
$3\overrightarrow{MA} - 3\overrightarrow{MB} = 3\overrightarrow{BA}$

\vspace{1em}

\noindent Donc $\| 5\overrightarrow{MG} \| = \| 3\overrightarrow{BA} \|$ \\
\phantom{Donc} $5MG = 3BA$ \\
\phantom{Donc} $5MG = 3 \times 5 \implies MG = \dfrac{3}{5}(5)$ \\
\phantom{Donc} \framebox{$MG = 3$}

\vspace{1.5em}

\noindent \textbf{L'ensemble des points $M$ est un cercle de centre $G$ et de rayon $3$.}

\vspace{2em}

\noindent \textbf{b)} $\| 2\overrightarrow{MA} + 3\overrightarrow{MB} \| = \| 3\overrightarrow{MA} + 2\overrightarrow{MB} \|$ \\
$2\overrightarrow{MA} + 3\overrightarrow{MB} = 5\overrightarrow{MG}$

\noindent Soit $H = \text{bar } \{(A, 3), (B, 2)\}$ \\
Alors $3 \overrightarrow{MA} + 2 \overrightarrow{MB} = (3 + 2) \overrightarrow{MH}$ \\
\phantom{Alors} $3 \overrightarrow{MA} + 2 \overrightarrow{MB} = 5 \overrightarrow{MH}$

\vspace{1em}

\noindent $\| 5 \overrightarrow{MG} \| = \| 5 \overrightarrow{MH} \|$ \\
$5 MG = 5 MH$ \\
$MG = MH$

\vspace{1.5em}

\noindent \textbf{Donc l'ensemble des points $M$ est la médiatrice de $[GH]$.}

\subsection*{Propriété 3 : Coordonnées du barycentre}

\noindent Soit $(O, \vec{i}, \vec{j})$ un repère orthonormé. \\
Si $G = \text{bar } \{(A, a), (B, b)\}$ alors :

\vspace{1.5em}

\noindent $x_G = \dfrac{a x_A + b x_B}{a + b}$

\vspace{1.5em}

\noindent $y_G = \dfrac{a y_A + b y_B}{a + b}$

\subsection*{\underline{\color{red}Preuve :}}

\noindent On a : $G = \text{bar } \{(A, a), (B, b)\} \iff \overrightarrow{OG} = \dfrac{1}{a+b} \left( a \overrightarrow{OA} + b \overrightarrow{OB} \right)$

\vspace{1em}

\noindent Or $\overrightarrow{OG} = \begin{pmatrix} x_G - x_O \\ y_G - y_O \end{pmatrix} \iff \overrightarrow{OG} = \begin{pmatrix} x_G \\ y_G \end{pmatrix}$

\vspace{1em}

\noindent $\overrightarrow{OA} = \begin{pmatrix} x_A - x_O \\ y_A - y_O \end{pmatrix} \iff \overrightarrow{OA} = \begin{pmatrix} x_A \\ y_A \end{pmatrix}$

\vspace{1em}

\noindent $\overrightarrow{OB} = \begin{pmatrix} x_B - x_O \\ y_B - y_O \end{pmatrix} \iff \overrightarrow{OB} = \begin{pmatrix} x_B \\ y_B \end{pmatrix}$

\vspace{1.5em}

\noindent Donc $\begin{pmatrix} x_G \\ y_G \end{pmatrix} = \dfrac{1}{a + b} \left( a \begin{pmatrix} x_A \\ y_A \end{pmatrix} + b \begin{pmatrix} x_B \\ y_B \end{pmatrix} \right)$

\vspace{1em}

\noindent \phantom{Donc $\begin{pmatrix} x_G \\ y_G \end{pmatrix}$} $= \dfrac{1}{a + b} \begin{pmatrix} a x_A + b x_B \\ a y_A + b y_B \end{pmatrix}$

\vspace{1em}

\noindent \phantom{Donc $\begin{pmatrix} x_G \\ y_G \end{pmatrix}$} $= \begin{pmatrix} \dfrac{a x_A + b x_B}{a + b} \\ \vspace{0.5em} \\ \dfrac{a y_A + b y_B}{a + b} \end{pmatrix}$

\noindent D'où $x_G = \dfrac{a x_A + b x_B}{a + b}$ \quad et \quad $y_G = \dfrac{a y_A + b y_B}{a + b}$

\vspace{2em}

\subsection*{Exemple :}

\noindent Soient $A \begin{pmatrix} 2 \\ 3 \end{pmatrix}$ et $B \begin{pmatrix} -2 \\ -1 \end{pmatrix}$ dans un repère orthonormé $(O, \vec{i}, \vec{j})$. \\
Soit $G$ barycentre des points pondérés $(A, 2)$ et $(B, 3)$. \\
Déterminer les coordonnées du point $G$.

\vspace{2em}

\begin{center}
    \textbf{Correction}
\end{center}

\noindent On a : $G = \text{bar } \{(A, 2), (B, 3)\}$ donc :

\vspace{1.5em}

\noindent $x_G = \dfrac{2 x_A + 3 x_B}{2 + 3} = \dfrac{2(2) + 3(-2)}{2 + 3} = -\dfrac{2}{5}$

\vspace{2em}

\noindent $y_G = \dfrac{2 y_A + 3 y_B}{2 + 3} = \dfrac{2(3) + 3(-1)}{2 + 3} = \dfrac{3}{5}$

\vspace{2em}

\noindent D'où \framebox{$G = \begin{pmatrix} -2/5 \\ 3/5 \end{pmatrix}$}

\section*{II- Barycentre de trois points pondérés}

\subsection*{Théorème - définition}

\noindent Soient $(A, a)$, $(B, b)$ et $(C, c)$ 3 points pondérés tel que $a + b + c \neq 0$. Alors il existe un unique point $G$ du plan tel que :
\begin{center}
    $a \overrightarrow{GA} + b \overrightarrow{GB} + c \overrightarrow{GC} = \vec{0}$
\end{center}

\noindent Le point $G$ est appelé barycentre des points $(A, a)$, $(B, b)$ et $(C, c)$, on note $G = \text{bar } \{(A, a), (B, b), (C, c)\}$ ou bien :
\begin{center}
    $G = \text{bar } \begin{array}{|c|c|c|}
        \hline
        A & B & C \\
        \hline
        a & b & c \\
        \hline
    \end{array}$
\end{center}

\vspace{1em}

\subsection*{Preuve : (TPE)}

\subsection*{Remarque :}

\noindent L'isobarycentre de 3 points non alignés est le centre de gravité du triangle formé par les points.

\vspace{1em}

\noindent \textbf{Ex :} \framebox{$ \overrightarrow{GA} + \overrightarrow{GB} + \overrightarrow{GC} = \vec{0} $}

\vspace{1em}

\noindent Donc $G = \text{bar } \{(A, 1), (B, 1), (C, 1)\}$ \\
$G$ est le centre de gravité du triangle $ABC$.

\subsection*{Propriété 1 : Réduction de la somme $a\overrightarrow{MA} + b\overrightarrow{MB} + c\overrightarrow{MC}$}

\noindent Soient $(A, a)$, $(B, b)$ et $(C, c)$ trois points pondérés tel que $a + b + c \neq 0$. \\
Pour tout point $M$ du plan on a : \\
$G = \text{bar } \{(A, a), (B, b), (C, c)\} \iff a\overrightarrow{MA} + b\overrightarrow{MB} + c\overrightarrow{MC} = (a + b + c)\overrightarrow{MG}$

\vspace{1em}

\noindent \textit{Preuve : (TPE)}

\vspace{1.5em}

\noindent \textbf{Remarque :} Compte tenu de la propriété 1, on a : \\
$G = \text{bar } \{(A, a), (B, b), (C, c)\} \iff \overrightarrow{MG} = \dfrac{1}{(a + b + c)} (a\overrightarrow{MA} + b\overrightarrow{MB} + c\overrightarrow{MC})$

\vspace{1.5em}

\noindent \textbf{Par conséquent :}

\vspace{1em}

\noindent $\overrightarrow{AG} = \dfrac{1}{(a + b + c)} (b\overrightarrow{AB} + c\overrightarrow{AC})$

\vspace{1.5em}

\noindent $\overrightarrow{BG} = \dfrac{1}{(a + b + c)} (a\overrightarrow{BA} + c\overrightarrow{BC})$

\vspace{1.5em}

\noindent $\overrightarrow{CG} = \dfrac{1}{(a + b + c)} (a\overrightarrow{CA} + b\overrightarrow{CB})$

\vspace{2em}

\noindent Pour la construction du barycentre $G$, souvent l'une de ces égalités est appliquée.

\subsection*{Ex :}

\noindent Soit $G = \text{bar } \{(A, 2), (B, 1), (C, 1)\}$ \\
Construire le point $G$.

\begin{center}
    \underline{Correction}
\end{center}

\noindent On a : $G = \text{bar } \{(A, 2), (B, 1), (C, 1)\}$ \\
donc $\overrightarrow{AG} = \dfrac{1}{2+1+1} (1 \overrightarrow{AB} + 1 \overrightarrow{AC})$ \\
\phantom{donc} $\overrightarrow{AG} = \dfrac{1}{4} (\overrightarrow{AB} + \overrightarrow{AC})$

\vspace{2em}

\subsection*{Propriété 2 : Coordonnées du barycentre}

\noindent Soit $G = \text{bar } \{(A, a), (B, b), (C, c)\}$ dans un repère orthonormé $(O, \vec{i}, \vec{j})$. Alors :

\vspace{1.5em}

\begin{center}
    \color{red}
    \framebox{\color{black} $x_G = \dfrac{ax_A + bx_B + cx_C}{a + b + c}$} 
    \quad et \quad 
    \framebox{\color{black} $y_G = \dfrac{ay_A + by_B + cy_C}{a + b + c}$}
\end{center}

\subsubsection*{Exemple :}

\noindent Soient $A \begin{pmatrix} 0 \\ 1 \end{pmatrix}$, $B \begin{pmatrix} 1 \\ 0 \end{pmatrix}$ et $C \begin{pmatrix} -1 \\ 2 \end{pmatrix}$ \\
Si $G = \text{bar } \{(A, 2), (B, 2), (C, 3)\}$ \\
Alors : $x_G = \dfrac{2x_A + 2x_B + 3x_C}{2 + 2 + 3}$ 

\vspace{1em}

\noindent et $y_G = \dfrac{2y_A + 2y_B + 3y_C}{2 + 2 + 3}$

\vspace{1.5em}

\noindent $x_G = \dfrac{2(0) + 2(1) + 3(-1)}{7} = -\dfrac{1}{7}$

\vspace{1.5em}

\noindent et $y_G = \dfrac{2(1) + 2(0) + 3(2)}{7} = \dfrac{8}{7}$

\vspace{1.5em}

\noindent D'où $G \begin{pmatrix} -1/7 \\ 8/7 \end{pmatrix}$

\vspace{2em}

\subsection*{Propriété 3 : Homogénéité}

\noindent Le barycentre de 3 points pondérés est inchangé lorsqu'on multiplie tous les points par un nombre réel non nul.

\vspace{1em}

\noindent $G = \text{bar } \{(A, a), (B, b), (C, c)\} \iff G = \text{bar } \{(A, ka), (B, kb), (C, kc)\}$

\subsubsection*{Exemple :}

\noindent $G = \text{bar } \{(A, -1), (B, 2), (C, 1)\}$. Alors : \\
$G = \text{bar } \{(A, -2), (B, 4), (C, 2)\}$ \quad ($k = 2$)

\vspace{2em}

\subsection*{Propriété 4 : Barycentre partiel (ou association du barycentre)}

\noindent Soient $a, b$ et $c$ des réels tels que : $a + b \neq 0$ et $a + b + c \neq 0$.

\vspace{1em}

\noindent Si $G = \text{bar } \{(A, a), (B, b), (C, c)\}$ et \\
$H = \text{bar } \{(A, a), (B, b)\}$ alors :

\vspace{1.5em}

\begin{center}
    \color{red}
    \framebox{\color{black} $G = \text{bar } \{(H, a + b), (C, c)\}$}
\end{center}

\vspace{2em}

\subsection*{Preuve :}

\noindent On a : $G = \text{bar } \{(A, a), (B, b), (C, c)\}$ \\
alors $a \overrightarrow{GA} + b \overrightarrow{GB} + c \overrightarrow{GC} = \vec{0}$ \\
$a (\overrightarrow{GH} + \overrightarrow{HA}) + b (\overrightarrow{GH} + \overrightarrow{HB}) + c \overrightarrow{GC} = \vec{0}$

\noindent $a \overrightarrow{GH} + a \overrightarrow{HA} + b \overrightarrow{GH} + b \overrightarrow{HB} + c \overrightarrow{GC} = \vec{0}$ \\
$(a + b) \overrightarrow{GH} + a \overrightarrow{HA} + b \overrightarrow{HB} + c \overrightarrow{GC} = \vec{0}$

\vspace{1em}

\noindent Sachant que $H = \text{bar } \{(A, a), (B, b)\}$ alors il en découle que $a \overrightarrow{HA} + b \overrightarrow{HB} = \vec{0}$

\vspace{1em}

\noindent Donc, par la suite on a : \\
$(a + b) \overrightarrow{GH} + \vec{0} + c \overrightarrow{GC} = \vec{0}$ \\
$(a + b) \overrightarrow{GH} + c \overrightarrow{GC} = \vec{0}$

\vspace{1em}

\noindent D'où $G = \text{bar } \{(H, a + b), (C, c)\}$

\vspace{2em}

\subsubsection*{Exemple :}

\noindent Si $G = \text{bar } \{(A, 2), (B, 1), (C, -1)\}$ \\
et $H = \text{bar } \{(A, 2), (B, 1)\}$ \\
alors $G = \text{bar } \{(H, 2 + 1), (C, -1)\}$

\vspace{1.5em}

\begin{center}
    \framebox{$G = \text{bar } \{(H, 3), (C, -1)\}$}
\end{center}

\section*{III - Utilisation de la propriété du barycentre partiel}

En général, pour montrer que 3 points sont alignés ou que 3 droites sont concourantes, on applique la propriété du barycentre partiel.

\vspace{1em}

\subsection*{1) 3 points alignés}

\noindent \textbf{N.B :} Pour montrer que 3 points sont alignés, il suffit de montrer que l'un des points est le barycentre des deux autres points en appliquant la propriété du barycentre partiel.

\vspace{2em}

\begin{center}
    \color{red} \Large \underline{Application}
\end{center}

\noindent Soit $ABC$ un triangle et $J$ le milieu de $[AC]$, $I$ et $K$ les points tels que $\overrightarrow{AI} = \dfrac{2}{3} \overrightarrow{AB}$ et $\overrightarrow{CK} = \dfrac{3}{4} \overrightarrow{CJ}$.

\vspace{1em}

\begin{itemize}
    \item[a)] Montrer que $J = \text{bar } \{(A, 1), (C, 1)\}$, $I = \text{bar } \{(A, 1), (B, 2)\}$ et $K = \text{bar } \{(I, 3), (C, 1)\}$.
    \item[b)] Montrer que $K = \text{bar } \{(A, 1), (B, 2), (C, 1)\}$.
    \item[c)] En déduire que $B, J$ et $K$ sont alignés.
\end{itemize}

\section*{Correction}

\noindent \textbf{1) *$_1$ Montrons que $I = \text{bar } \{(A, 1), (C, 1)\}$} \\
On a : $I$ est le milieu de $[AC]$ donc $I$ est l'isobarycentre des points $A$ et $C$. \\
D'où \framebox{$I = \text{bar } \{(A, 1), (C, 1)\}$}

\vspace{2em}

\noindent \textbf{*$_2$ Montrons que $J = \text{bar } \{(A, 1), (B, 2)\}$}

\vspace{1em}

\noindent \underline{Méthode 1 :} \\
On a : $\overrightarrow{AJ} = \dfrac{2}{3} \overrightarrow{AB} \iff 3 \overrightarrow{AJ} = 2 \overrightarrow{AB}$ \\
\phantom{On a : $\overrightarrow{AJ} = \dfrac{2}{3} \overrightarrow{AB} \iff$} $3 \overrightarrow{AJ} = 2 (\overrightarrow{AJ} + \overrightarrow{JB})$ \\
\phantom{On a : $\overrightarrow{AJ} = \dfrac{2}{3} \overrightarrow{AB} \iff$} $3 \overrightarrow{AJ} = 2 \overrightarrow{AJ} + 2 \overrightarrow{JB}$ \\
\phantom{On a : $\overrightarrow{AJ} = \dfrac{2}{3} \overrightarrow{AB} \iff$} $3 \overrightarrow{AJ} - 2 \overrightarrow{AJ} - 2 \overrightarrow{JB} = \vec{0}$ \\
\phantom{On a : $\overrightarrow{AJ} = \dfrac{2}{3} \overrightarrow{AB} \iff$} $\overrightarrow{AJ} - 2 \overrightarrow{JB} = \vec{0}$ \\
\phantom{On a : $\overrightarrow{AJ} = \dfrac{2}{3} \overrightarrow{AB} \iff$} $-\overrightarrow{JA} - 2 \overrightarrow{JB} = \vec{0}$ \quad or $-1 - 2 \neq 0$

\vspace{1em}

\noindent Donc $J = \text{bar } \{(A, -1), (B, -2)\}$ \\
D'où par l'homogénéité du barycentre, on a :
\begin{center}
    \framebox{$J = \text{bar } \{(A, 1), (B, 2)\}$}
\end{center}

\noindent \underline{Méthode 2 :}

\noindent \textbf{Rappel :} \framebox{$G = \text{bar } \{(A, a), (B, b)\} \iff \overrightarrow{AG} = \dfrac{b}{a + b} \overrightarrow{AB}$}

\vspace{1.5em}

\noindent On a : $\overrightarrow{AJ} = \dfrac{2}{3} \overrightarrow{AB} \iff \overrightarrow{AJ} = \dfrac{2}{1 + 2} \overrightarrow{AB}$

\vspace{1em}

\noindent D'où \framebox{$J = \text{bar } \{(A, 1), (B, 2)\}$}

\vspace{2em}

\noindent \textbf{*$_3$ Montrons que $K = \text{bar } \{(J, 3), (C, 1)\}$}

\noindent \underline{Méthode 1 :} Montrer que $3 \overrightarrow{JK} + \overrightarrow{JC} = \vec{0}$ (TPE)

\noindent \underline{Méthode 2 :}

\noindent On a : $\overrightarrow{CK} = \dfrac{3}{4} \overrightarrow{CJ} \iff \overrightarrow{CK} = \dfrac{3}{1 + 3} \overrightarrow{CJ}$

\vspace{1.5em}

\noindent Donc $K = \text{bar } \{(C, 1), (J, 3)\}$ \\
D'où \framebox{$K = \text{bar } \{(J, 3), (C, 1)\}$}

\vspace{2em}

\noindent \textbf{b) Montrons que $K = \text{bar } \{(A, 1), (B, 2), (C, 1)\}$}

\noindent On a : $K = \text{bar } \{(J, 3), (C, 1)\}$ \\
\phantom{On a : } et $J = \text{bar } \{(A, 1), (B, 2)\}$. Alors d'après la propriété \\
du barycentre partiel on a : \\
$K = \text{bar } \{(A, 1), (B, 2), (C, 1)\}$

\vspace{2em}

\noindent \textbf{c) Déduisons-en que les points $B, I$ et $K$ sont alignés.}

\vspace{1em}

\noindent On a : $K = \text{bar } \{(A, 1), (B, 2), (C, 1)\}$ \\
\phantom{On a : } $K = \text{bar } \{(A, 1), (C, 1), (B, 2)\}$ \\
\phantom{On a : } or $I = \text{bar } \{(A, 1), (C, 1)\}$

\vspace{1em}

\noindent D'après la propriété du barycentre partiel on a : \\
$K = \text{bar } \{(I, 1 + 1), (B, 2)\}$ \\
$K = \text{bar } \{(I, 2), (B, 2)\}$

\vspace{1.5em}

\noindent \textbf{D'où $B, I$ et $K$ sont alignés.}

\subsection*{2) Droites concourantes}

\noindent \textbf{N.B :} Pour montrer que des droites sont concourantes, il suffit de trouver un point commun à toutes ces droites. En général, pour le prouver, on applique la propriété du barycentre partiel.

\vspace{2em}

\begin{center}
    \color{red} \Large \underline{Application}
\end{center}

\noindent Soit $ABC$ un triangle, $I, J, K$ des points tels que : \\
$\overrightarrow{AI} = \dfrac{1}{3} \overrightarrow{AB}$ , $\overrightarrow{AJ} = \dfrac{3}{4} \overrightarrow{AC}$ et $\overrightarrow{BK} = \dfrac{6}{7} \overrightarrow{BC}$

\vspace{1em}

\begin{itemize}
    \item[1)] Montrer que $I = \text{bar } \{(A, 2), (B, 1)\}$, \\
    $J = \text{bar } \{(A, 2), (C, 6)\}$ et $K = \text{bar } \{(B, 1), (C, 6)\}$
    \item[2)] On pose $G = \text{bar } \{(A, 2), (B, 1), (C, 6)\}$
    \begin{itemize}
        \item[a)] Montrer que $G \in (AK)$, $G \in (BJ)$ et $G \in (IC)$
        \item[b)] Justifier que les droites $(AK)$, $(BJ)$ et $(IC)$ sont concourantes.
    \end{itemize}
\end{itemize}

\section*{Correction}

\noindent \textbf{1) Montrons les relations de barycentre :}

\vspace{0.5em}
\noindent $\bullet$ \textbf{Pour $I$ :} On a $\overrightarrow{AI} = \frac{1}{3} \overrightarrow{AB} \iff 3\overrightarrow{AI} = \overrightarrow{AB} \iff 3\overrightarrow{AI} = \overrightarrow{AI} + \overrightarrow{IB}$ \\
$\iff 2\overrightarrow{AI} - \overrightarrow{IB} = \vec{0} \iff -2\overrightarrow{IA} - \overrightarrow{IB} = \vec{0} \iff 2\overrightarrow{IA} + \overrightarrow{IB} = \vec{0}$ \\
D'où \framebox{$I = \text{bar } \{(A, 2), (B, 1)\}$}

\vspace{1em}
\noindent $\bullet$ \textbf{Pour $J$ :} On a $\overrightarrow{AJ} = \frac{3}{4} \overrightarrow{AC} \iff 4\overrightarrow{AJ} = 3\overrightarrow{AC} \iff 4\overrightarrow{AJ} = 3(\overrightarrow{AJ} + \overrightarrow{JC})$ \\
$\iff \overrightarrow{AJ} - 3\overrightarrow{JC} = \vec{0} \iff -\overrightarrow{JA} - 3\overrightarrow{JC} = \vec{0} \iff \overrightarrow{JA} + 3\overrightarrow{JC} = \vec{0}$ \\
Par homogénéité ($k=2$) : \framebox{$J = \text{bar } \{(A, 2), (C, 6)\}$}

\vspace{1em}
\noindent $\bullet$ \textbf{Pour $K$ :} On a $\overrightarrow{BK} = \frac{6}{7} \overrightarrow{BC} \iff 7\overrightarrow{BK} = 6(\overrightarrow{BK} + \overrightarrow{KC})$ \\
$\iff \overrightarrow{BK} - 6\overrightarrow{KC} = \vec{0} \iff -\overrightarrow{KB} - 6\overrightarrow{KC} = \vec{0} \iff \overrightarrow{KB} + 6\overrightarrow{KC} = \vec{0}$ \\
D'où \framebox{$K = \text{bar } \{(B, 1), (C, 6)\}$}

\vspace{2em}

\noindent \textbf{2) a) Montrons l'appartenance de $G$ aux droites :}

\vspace{0.5em}
\noindent $G = \text{bar } \{(A, 2), (B, 1), (C, 6)\}$
\begin{itemize}
    \item $G = \text{bar } \{(A, 2), \underbrace{(B, 1), (C, 6)}_{K, 7}\} \implies G = \text{bar } \{(A, 2), (K, 7)\} \implies \mathbf{G \in (AK)}$
    \item $G = \text{bar } \{(B, 1), \underbrace{(A, 2), (C, 6)}_{J, 8}\} \implies G = \text{bar } \{(B, 1), (J, 8)\} \implies \mathbf{G \in (BJ)}$
    \item $G = \text{bar } \{(C, 6), \underbrace{(A, 2), (B, 1)}_{I, 3}\} \implies G = \text{bar } \{(C, 6), (I, 3)\} \implies \mathbf{G \in (IC)}$
\end{itemize}

\vspace{1em}
\noindent \textbf{b) Justification :} \\
Le point $G$ appartient simultanément aux trois droites $(AK)$, $(BJ)$ et $(IC)$. Par conséquent, ces trois droites sont concourantes en $G$.

\end{document}